\chapter{Реализация микросервисного проекта РГР}

\section{Постановка задачи}

В расчетно-графической работе разработан проект \texttt{likes-s18}. Основная идея проекта --
разделить функциональность социальной платформы на несколько независимых сервисов и связать их
через разные протоколы взаимодействия. В результате получилась система, в которой каждый сервис
отвечает за отдельный домен и может развиваться независимо.

В состав проекта входят три сервиса:
\begin{itemize}
    \item \texttt{user-service} -- регистрация и получение пользователей;
    \item \texttt{post-service} -- создание, получение и удаление постов;
    \item \texttt{like-service} -- создание лайков, получение лайков и подсчет их количества.
\end{itemize}

Ключевой пользовательский сценарий выглядит следующим образом: клиент создает пост, затем
создает лайки к этому посту, после чего запрашивает пост по идентификатору. При выдаче поста
\texttt{post-service} обращается к \texttt{like-service} по gRPC и добавляет в ответ поле
\texttt{likes\_count}. Благодаря этому клиенту не нужно выполнять дополнительный запрос для
подсчета лайков.

\section{Архитектура системы}

Разделение на сервисы выполнено по бизнес-доменам. Пользователи, посты и лайки имеют разную
ответственность, разные модели данных и потенциально разную нагрузку. Поэтому каждый домен
вынесен в отдельный сервис и работает со своей базой PostgreSQL \cite{postgres_docs}.

Архитектура проекта построена по следующей схеме:
\begin{enumerate}
    \item внешний клиент обращается к REST API сервисов;
    \item \texttt{post-service} хранит и возвращает данные постов;
    \item \texttt{like-service} хранит данные лайков;
    \item при получении конкретного поста \texttt{post-service} выполняет gRPC-вызов в
    \texttt{like-service};
    \item \texttt{like-service} подсчитывает лайки в своей базе данных и возвращает результат.
\end{enumerate}

Такое разделение снижает связанность. Например, изменение логики подсчета лайков не требует
изменять код сервиса пользователей. Аналогично, если нужно масштабировать обработку лайков, можно
увеличить количество экземпляров только \texttt{like-service}, не затрагивая остальные части
системы \cite{week17_arch}.

\section{Выбор протоколов взаимодействия}

Для внешнего API выбран REST, так как он прост в ручном тестировании, хорошо поддерживается
браузерами, HTTP-клиентами и FastAPI \cite{fastapi_docs}. Для внутреннего взаимодействия между
\texttt{post-service} и \texttt{like-service} выбран gRPC. Он удобен тем, что контракт описывается
в \texttt{.proto}-файле, а клиентский и серверный код используют сгенерированные классы
\cite{grpc_docs}.

Сравнение используемых протоколов приведено в таблице \ref{tab:rest-grpc}.

\begin{table}[h!]
    \caption{Сравнение REST и gRPC в проекте}
    \label{tab:rest-grpc}
    \noindent
    \begin{tabular*}{\textwidth}{|p{0.22\textwidth}|p{0.30\textwidth}|p{0.34\textwidth}|}
        \hline
        Критерий & REST & gRPC \\
        \hline
        Назначение & Внешнее API & Внутренний вызов подсчета лайков \\
        \hline
        Формат данных & JSON & Protobuf \\
        \hline
        Контракт & HTTP endpoint-ы & \texttt{likes.proto} \\
        \hline
        Удобство & Просто проверять через curl & Удобно для сервис-сервис связи \\
        \hline
    \end{tabular*}
\end{table}

\section{REST API сервисов}

REST API реализовано на FastAPI. Каждый сервис создает собственное приложение и описывает набор
endpoint-ов для работы со своими сущностями. Например, в \texttt{like-service} реализованы
методы создания лайка, получения списка лайков, получения лайка по идентификатору, удаления лайка
и отдельный HTTP endpoint для подсчета лайков по посту.

Фрагмент REST-обработчика создания лайка приведен в листинге \ref{lst:create-like}.

\begin{lstlisting}[language=Python, caption={Создание лайка в like-service}, label={lst:create-like}]
@app.post("/likes")
def create_like(post_id: int, db: Session = Depends(get_db)):
    new_like = Like(post_id=post_id)
    db.add(new_like)
    db.commit()
    db.refresh(new_like)
    return new_like
\end{lstlisting}

В этом обработчике FastAPI получает параметр \texttt{post\_id}, открывает сессию SQLAlchemy
через зависимость \texttt{Depends(get\_db)}, создает объект модели \texttt{Like}, сохраняет его
в базе данных и возвращает созданную запись. Аналогичный подход используется в сервисах
пользователей и постов \cite{sqlalchemy_docs}.

\section{gRPC-контракт и межсервисный вызов}

gRPC-контракт описывает один метод \texttt{GetLikesCount}. На вход метод принимает идентификатор
поста, а в ответ возвращает количество лайков. Контракт приведен в листинге
\ref{lst:proto-contract}.

\begin{lstlisting}[caption={gRPC-контракт likes.proto}, label={lst:proto-contract}]
syntax = "proto3";

service LikeService {
  rpc GetLikesCount (LikeRequest) returns (LikeResponse);
}

message LikeRequest {
  int32 post_id = 1;
}

message LikeResponse {
  int32 count = 1;
}
\end{lstlisting}

Серверная часть gRPC находится в \texttt{like-service}. Она получает \texttt{post\_id}, выполняет
запрос к таблице лайков и возвращает количество найденных записей. Основной фрагмент реализации
показан в листинге \ref{lst:grpc-server}.

\begin{lstlisting}[language=Python, caption={Метод GetLikesCount на стороне like-service}, label={lst:grpc-server}]
def GetLikesCount(self, request, context):
    db = SessionLocal()
    count = db.query(Like).filter(
        Like.post_id == request.post_id
    ).count()
    return likes_pb2.LikeResponse(count=count)
\end{lstlisting}

На стороне \texttt{post-service} создается gRPC-канал до контейнера \texttt{like-service}.
Затем вызывается метод \texttt{GetLikesCount}, а результат добавляется в REST-ответ клиенту.
Этот фрагмент приведен в листинге \ref{lst:grpc-call}.

\begin{lstlisting}[language=Python, caption={Вызов gRPC в post-service}, label={lst:grpc-call}]
channel = grpc.insecure_channel("like-service:50051")
stub = likes_pb2_grpc.LikeServiceStub(channel)
response = stub.GetLikesCount(
    likes_pb2.LikeRequest(post_id=post_id)
)
likes_count = response.count
\end{lstlisting}

Важная деталь состоит в том, что имя \texttt{like-service} используется как сетевое имя внутри
Docker Compose. Благодаря встроенной DNS-сети Compose сервисы могут обращаться друг к другу по
именам, указанным в файле \texttt{docker-compose.yml}. В Kubernetes используется тот же подход:
внутри namespace сервис доступен по DNS-имени \texttt{like-service}.

\section{Запуск REST и gRPC внутри like-service}

\texttt{like-service} должен одновременно обслуживать HTTP-запросы FastAPI и gRPC-запросы от
\texttt{post-service}. Для этого gRPC-сервер запускается в отдельном потоке во время старта
FastAPI-приложения. Фрагмент запуска показан в листинге \ref{lst:lifespan}.

\begin{lstlisting}[language=Python, caption={Запуск gRPC-сервера в lifespan FastAPI}, label={lst:lifespan}]
@asynccontextmanager
async def lifespan(app: FastAPI):
    threading.Thread(
        target=start_grpc_server,
        daemon=True
    ).start()
    print("gRPC thread started")
    yield
    print("Application shutdown")
\end{lstlisting}

Такое решение позволяет поднять оба интерфейса в одном контейнере: REST API работает на порту
\texttt{8000}, а gRPC-сервер слушает порт \texttt{50051}. При этом внешнему клиенту доступен
REST-интерфейс, а внутренние сервисы используют gRPC.

\section{Отказоустойчивость и healthcheck}

В \texttt{post-service} gRPC-вызов обернут в \texttt{try/except}. Если \texttt{like-service}
временно недоступен, сервис постов не падает, а возвращает пост с \texttt{likes\_count = 0}.
Это пример деградации функциональности: клиент все равно получает основную информацию о посте,
но без точного количества лайков.

\begin{lstlisting}[language=Python, caption={Обработка ошибки gRPC-вызова}, label={lst:grpc-fallback}]
try:
    response = stub.GetLikesCount(
        likes_pb2.LikeRequest(post_id=post_id)
    )
    likes_count = response.count
except Exception as e:
    print("gRPC ERROR:", e)
    likes_count = 0
\end{lstlisting}

Для проверки состояния сервисов реализован endpoint \texttt{/health}. Он не просто возвращает
статический ответ, а выполняет запрос \texttt{SELECT 1} к базе данных. Поэтому healthcheck
проверяет не только доступность HTTP-приложения, но и подключение к хранилищу.

\begin{lstlisting}[language=Python, caption={Проверка состояния сервиса}, label={lst:health}]
@app.get("/health")
def health(db: Session = Depends(get_db)):
    try:
        db.execute(text("SELECT 1"))
    except Exception as e:
        raise HTTPException(
            status_code=503,
            detail=f"database is unavailable: {e}"
        )
    return {"status": "ok", "service": "post-service"}
\end{lstlisting}

\section{Контейнеризация через Docker}

Каждый сервис упакован в отдельный Docker-образ. Dockerfile задает базовый Python-образ,
рабочую директорию, установку зависимостей, копирование исходного кода и команду запуска
FastAPI-приложения через \texttt{uvicorn}. Пример Dockerfile для \texttt{post-service} приведен
в листинге \ref{lst:dockerfile}.

\begin{lstlisting}[caption={Dockerfile сервиса постов}, label={lst:dockerfile}]
FROM python:3.12

WORKDIR /app

COPY ./post_service/requirements.txt ./requirements.txt

RUN pip install --no-cache-dir -r requirements.txt

COPY ./post_service /app
COPY ./proto /app/proto

CMD ["python", "-m", "uvicorn", "posts:app",
     "--host", "0.0.0.0", "--port", "8000"]
\end{lstlisting}

Отдельная упаковка сервисов дает несколько преимуществ: окружение становится воспроизводимым,
зависимости не смешиваются между сервисами, а запуск приложения не зависит от локальной настройки
Python на машине разработчика \cite{docker_compose_docs}.

\section{Оркестрация через Docker Compose}

Для локального запуска всех компонентов используется Docker Compose. В одном файле описаны три
API-сервиса и три базы данных PostgreSQL. Внешние порты проброшены следующим образом:
\begin{itemize}
    \item \texttt{like-service}: REST -- \texttt{8001}, gRPC -- \texttt{50051};
    \item \texttt{post-service}: REST -- \texttt{8002};
    \item \texttt{user-service}: REST -- \texttt{8003}.
\end{itemize}

Фрагмент конфигурации \texttt{like-service} приведен в листинге \ref{lst:compose-like}.

\begin{lstlisting}[caption={Фрагмент docker-compose.yml для like-service}, label={lst:compose-like}]
like-service:
  build:
    context: .
    dockerfile: like_service/Dockerfile
  ports:
    - "8001:8000"
    - "50051:50051"
  depends_on:
    like-db:
      condition: service_healthy
  healthcheck:
    test: ["CMD-SHELL", "python -c \"import urllib.request; urllib.request.urlopen('http://127.0.0.1:8000/health', timeout=3).read()\""]
    interval: 10s
    timeout: 5s
    retries: 5
    start_period: 20s
\end{lstlisting}

Параметр \texttt{depends\_on} с условием \texttt{service\_healthy} нужен для корректного порядка
запуска. API-контейнер стартует только после того, как контейнер базы данных пройдет проверку
готовности. Для PostgreSQL используется команда \texttt{pg\_isready}; она позволяет понять, что
СУБД уже принимает подключения.

Запуск всей системы выполняется одной командой из каталога приложения:
\begin{lstlisting}[caption={Запуск контейнеров проекта}]
docker compose up -d --build
\end{lstlisting}

\section{Настройка CI}

Для проекта настроен GitHub Actions workflow, который автоматически проверяет РГР при изменении
кода сервисов, вариантов, зависимостей или самого файла CI. Pipeline выполняется в окружении
\texttt{ubuntu-latest} и использует переменные \texttt{GROUP=433} и \texttt{STUDENT\_ID=s18},
чтобы тесты брали правильный вариант проекта.

Основные шаги CI:
\begin{enumerate}
    \item checkout репозитория;
    \item установка Python 3.12;
    \item установка зависимостей из \texttt{requirements.txt};
    \item запуск тестов командой \texttt{make test WEEK=17};
    \item сборка и запуск контейнеров через Docker Compose;
    \item smoke-тест REST + gRPC-сценария;
    \item вывод логов при ошибке;
    \item остановка контейнеров и удаление volume-ов.
\end{enumerate}

Фрагмент настройки CI приведен в листинге \ref{lst:ci}.

\begin{lstlisting}[caption={Фрагмент GitHub Actions CI}, label={lst:ci}]
- name: Run tests
  run: make test WEEK=17

- name: Build and start services
  run: docker compose -f "$APP_DIR/docker-compose.yml" up --build -d

- name: Smoke test REST + gRPC path
  run: |
    sleep 20
    docker compose -f "$APP_DIR/docker-compose.yml" exec -T like-service python -c "import socket; s=socket.socket(); s.settimeout(3); assert s.connect_ex(('127.0.0.1',50051))==0"
    docker compose -f "$APP_DIR/docker-compose.yml" exec -T post-service python -c "import requests; p=requests.post('http://127.0.0.1:8000/posts', params={'title':'ci-smoke','content':'ok'}).json(); r=requests.get(f\"http://127.0.0.1:8000/posts/{p['id']}\"); assert r.status_code==200; assert 'likes_count' in r.json()"
\end{lstlisting}

Smoke-тест важен тем, что он проверяет не только наличие файлов или прохождение статического
теста, но и фактический запуск системы в контейнерах. Сначала проверяется, что gRPC-порт
\texttt{50051} открыт внутри \texttt{like-service}. Затем внутри \texttt{post-service} создается
пост и выполняется запрос на получение этого поста. Успешная проверка поля \texttt{likes\_count}
подтверждает, что цепочка \texttt{post-service -> gRPC -> like-service} работает.

\section{Kubernetes-манифесты}

Для кластерного развертывания добавлена директория \texttt{weeks/week-17/k8s}, в которой
описаны все ресурсы проекта:
\begin{itemize}
    \item namespace \texttt{likes-s18};
    \item Deployment + Service для \texttt{user-service}, \texttt{post-service}, \texttt{like-service};
    \item Deployment + Service для \texttt{user-db}, \texttt{post-db}, \texttt{like-db};
    \item файл \texttt{kustomization.yaml} для применения полного набора манифестов.
\end{itemize}

Фрагмент манифеста \texttt{like-service} приведен в листинге \ref{lst:k8s-like-service}.

\begin{lstlisting}[caption={Фрагмент Kubernetes-манифеста like-service}, label={lst:k8s-like-service}]
apiVersion: apps/v1
kind: Deployment
metadata:
  name: like-service
  namespace: likes-s18
spec:
  replicas: 2
  template:
    spec:
      containers:
        - name: like-service
          image: ghcr.io/example/likes-s18-like-service:latest
          ports:
            - containerPort: 8000
            - containerPort: 50051
          readinessProbe:
            httpGet:
              path: /health
              port: 8000
\end{lstlisting}

В этом манифесте фиксируется желаемое число реплик, образ контейнера, открытые порты и
проверка готовности через endpoint \texttt{/health}. Наличие readiness/liveness-проб позволяет
оркестратору не отправлять трафик в неготовые pod-ы и автоматически перезапускать их при сбоях
\cite{kubernetes_docs}.

Базы данных также описаны как отдельные Deployment/Service-ресурсы. Для PostgreSQL используется
probe на основе \texttt{pg\_isready}, что обеспечивает корректный запуск API только после
готовности хранилища.

\section{Continuous Delivery (CD)}

Отдельный workflow \texttt{week-17-cd.yml} реализует конвейер доставки:
\begin{enumerate}
    \item запуск тестов;
    \item сборка и публикация трех образов в GHCR;
    \item применение Kubernetes-манифестов;
    \item обновление образов в Deployment через \texttt{kubectl set image};
    \item ожидание завершения rolling update через \texttt{kubectl rollout status}.
\end{enumerate}

Ключевой фрагмент CD приведен в листинге \ref{lst:cd-deploy}.

\begin{lstlisting}[caption={Шаги deploy в GitHub Actions CD}, label={lst:cd-deploy}]
- name: Apply manifests
  run: |
    kubectl apply -f weeks/week-17/k8s/namespace.yaml
    kubectl apply -f weeks/week-17/k8s/

- name: Update images to current commit
  run: |
    kubectl -n "${K8S_NAMESPACE}" set image deployment/like-service \
      like-service=ghcr.io/${{ github.repository }}/week-17-like-service:${{ github.sha }}

- name: Wait for rollout
  run: kubectl -n "${K8S_NAMESPACE}" rollout status deployment/like-service --timeout=180s
\end{lstlisting}

Таким образом, при изменении кода в ветке \texttt{main} pipeline доводит систему до нового
согласованного состояния в кластере: тесты проходят, образы публикуются, а деплой обновляется без
полного простоя за счет rolling update \cite{github_actions_docs,kubectl_docs}.

\section{Результаты выполнения РГР}

В результате выполнения РГР получена микросервисная система, которая:
\begin{itemize}
    \item реализует REST API для пользователей, постов и лайков;
    \item использует gRPC для внутреннего подсчета лайков;
    \item хранит данные каждого сервиса в отдельной базе PostgreSQL;
    \item запускается локально через Docker Compose;
    \item имеет healthcheck-и для API и баз данных;
    \item автоматически проверяется в GitHub Actions CI;
    \item разворачивается в Kubernetes через набор манифестов;
    \item поддерживает CD-сценарий публикации образов и rolling update.
\end{itemize}

Главным результатом является не только работоспособность отдельных endpoint-ов, но и проверенная
интеграция нескольких компонентов: HTTP API, gRPC, баз данных, контейнеров и CI-сценария.

\section{Ответы на вопросы самопроверки}

\textbf{1. Критерии разделения на сервисы.}
Разделение выполнено по бизнес-возможностям системы: пользователи, посты, лайки. Каждый сервис
владеет своим API и своей базой данных. Такая модель уменьшает связанность модулей и позволяет
масштабировать наиболее нагруженные части системы отдельно.

\textbf{2. Причины выбора стека.}
Python и FastAPI выбраны за высокую скорость разработки и встроенную поддержку OpenAPI.
PostgreSQL использован как стабильная реляционная СУБД. REST применен для клиентского API, а gRPC
-- для внутреннего вызова с явным контрактом через \texttt{.proto}.

\textbf{3. Обработка ошибок и устойчивость к сбоям.}
Критичный внешний endpoint сервиса постов реализует безопасное поведение при ошибке gRPC-вызова:
исключение перехватывается, сервис не падает, а клиент получает валидный ответ с нулевым
значением лайков.

\textbf{4. Подход к обновлению без простоя.}
В проект добавлен CD workflow для Kubernetes: после тестов публикуются новые образы, затем
выполняется \texttt{kubectl set image} и контролируется \texttt{kubectl rollout status}. За счет
readinessProbe и механизма rolling update трафик переключается на новые pod-ы постепенно, что
позволяет обновлять сервисы без полной остановки системы.

\textbf{5. Наиболее сложный этап интеграции.}
Сложность представлял одновременный запуск REST и gRPC в \texttt{like-service}. Для корректной
работы gRPC-сервер был вынесен в отдельный поток и запущен в lifecycle FastAPI-приложения.

\textbf{6. Возможные улучшения.}
При дальнейшем развитии проекта рекомендуется добавить централизованное логирование, сбор метрик
Prometheus/Grafana, retry/backoff и circuit breaker для внутренних вызовов, а также API Gateway
для маршрутизации и единой точки авторизации.
